\documentclass[../DoAn.tex]{subfiles}
\begin{document}

\section{Tự động chuyển hóa video bài giảng thành bài học tương tác}
\subsection{Đặt vấn đề}
Trong mô hình học trực tuyến thông thường, video bài giảng thường chỉ đóng vai trò là tài liệu truyền tải một chiều. Giảng viên có thể tải video lên, học viên có thể xem lại nhiều lần, nhưng hệ thống khó bảo đảm rằng học viên thực sự theo dõi và hiểu nội dung tại từng thời điểm của bài học. Nếu muốn biến một video thành bài học có tương tác, giảng viên thường phải thực hiện nhiều thao tác thủ công: nghe lại video, ghi phụ đề hoặc tóm tắt nội dung, xác định các đoạn kiến thức quan trọng, soạn câu hỏi cho từng đoạn và đặt câu hỏi vào đúng mốc thời gian. Quy trình này tốn thời gian, dễ bỏ sót ý chính và khó mở rộng khi số lượng bài giảng lớn.

Đối với hệ thống Dyadia, yêu cầu đặt ra không chỉ là lưu trữ và phát video mà còn phải biến video thành một bài học có khả năng kiểm tra, phản hồi và theo dõi mức độ tiếp thu của học viên. Bài toán khó nằm ở chỗ dữ liệu đầu vào chỉ là một tệp video, trong khi đầu ra cần có nhiều thành phần có cấu trúc: bản chép lời, các phân đoạn nội dung, hoạt động tương tác gắn với từng đoạn, đáp án, giải thích và cấu hình thời điểm xuất hiện. Nếu chỉ dùng một lời gọi mô hình ngôn ngữ lớn để sinh toàn bộ kết quả, hệ thống dễ gặp lỗi định dạng, khó kiểm soát chất lượng và khó cho phép giảng viên biên tập từng phần.

\subsection{Giải pháp}
Giải pháp được xây dựng là một quy trình xử lý nhiều giai đoạn, trong đó video bài giảng đi qua ba bước chính: phiên âm lời thoại, phân đoạn nội dung và sinh hoạt động tương tác. Mỗi bước có trách nhiệm rõ ràng và tạo ra một dạng dữ liệu trung gian có thể lưu trữ, xem lại và chỉnh sửa. Thay vì xem trí tuệ nhân tạo như một hộp đen tạo ra toàn bộ bài học, hệ thống tách quy trình thành các bước nhỏ để dễ kiểm soát và dễ phục hồi khi có lỗi.

Ở bước thứ nhất, hệ thống trích xuất lời thoại từ video bằng mô hình nhận dạng tiếng nói. Kết quả phiên âm được lưu lại như một nguồn dữ liệu độc lập, phục vụ cả việc hiển thị phụ đề và làm đầu vào cho các bước sau. Ở bước thứ hai, bản chép lời được phân đoạn theo ý nghĩa nội dung. Mỗi phân đoạn có mốc bắt đầu, mốc kết thúc và phần tóm tắt ngắn, giúp liên kết nội dung bài giảng với dòng thời gian của video. Ở bước thứ ba, mô hình ngôn ngữ lớn sinh các hoạt động tương tác cho từng phân đoạn dựa trên cấu hình của bài học, chẳng hạn số lượng câu hỏi, loại tương tác, độ khó và trọng tâm mà giảng viên mong muốn.

Một điểm quan trọng của giải pháp là hệ thống không khóa giảng viên vào kết quả sinh tự động. Sau khi quy trình hoàn thành, giảng viên có thể xem trước bản chép lời, các phân đoạn và từng hoạt động tương tác; từ đó chỉnh sửa nội dung, thay đổi đáp án, thêm hoạt động thủ công hoặc yêu cầu sinh lại cho một đoạn cụ thể. Cách làm này kết hợp ưu điểm của tự động hóa với vai trò kiểm duyệt chuyên môn của giảng viên. Trí tuệ nhân tạo đảm nhận phần tạo nháp tốn công, còn giảng viên giữ quyền quyết định cuối cùng trước khi xuất bản bài học.

Về mặt triển khai, các kết quả trung gian được lưu tách biệt theo bản chất dữ liệu. Dữ liệu nghiệp vụ như bài giảng, hoạt động tương tác, cấu hình, trạng thái phân tích được lưu trong PostgreSQL. Các dữ liệu lớn hoặc có cấu trúc thay đổi như bản chép lời và nội dung phân đoạn được lưu trong FoundationDB. Tệp video và âm thanh được lưu trong SeaweedFS thông qua giao thức tương thích S3. Cách phân tách này giúp hệ thống không ép mọi loại dữ liệu vào một cơ sở dữ liệu quan hệ duy nhất, đồng thời vẫn giữ được liên kết giữa các đối tượng thông qua định danh bài giảng, phân đoạn và hoạt động.

\subsection{Kết quả đạt được}
Kết quả của giải pháp là hệ thống đã hỗ trợ luồng nghiệp vụ hoàn chỉnh: giảng viên tạo bài giảng từ video, kích hoạt phân tích, theo dõi trạng thái xử lý, xem lại kết quả sinh ra, biên tập hoạt động và xuất bản bài học. Từ góc nhìn người dùng, một video thô được chuyển thành bài học có phụ đề, có phân đoạn nội dung và có các điểm tương tác xuất hiện trong quá trình học.

Đóng góp chính của phần này không nằm ở việc gọi riêng lẻ một mô hình trí tuệ nhân tạo, mà nằm ở cách tổ chức toàn bộ quy trình thành một pipeline có dữ liệu trung gian rõ ràng và có điểm kiểm soát cho giảng viên. Nhờ vậy, hệ thống có thể vận hành ổn định hơn, dễ kiểm thử hơn và phù hợp hơn với bối cảnh giáo dục, nơi nội dung sinh tự động cần được kiểm duyệt trước khi đưa tới học viên.

\section{Mô hình hoạt động tương tác mở rộng được}
\subsection{Đặt vấn đề}
Một hệ thống học bằng video chỉ có giá trị tương tác cao nếu nó hỗ trợ nhiều dạng hoạt động khác nhau. Trắc nghiệm một đáp án là dạng phổ biến và dễ chấm điểm, nhưng không đủ để đánh giá toàn diện quá trình học. Với nội dung ngoại ngữ hoặc kỹ năng mềm, hệ thống cần bài nghe, bài đọc, câu hỏi mở hoặc bài viết ngắn. Với nội dung kỹ thuật, hệ thống cần dạng điền khuyết, chọn nhiều đáp án, hoặc câu hỏi yêu cầu giải thích. Nếu mỗi loại hoạt động được cài đặt như một trường hợp đặc biệt rải rác trong cả máy chủ và giao diện, hệ thống sẽ nhanh chóng khó bảo trì, khó thêm loại mới và dễ phát sinh lỗi không thống nhất giữa các màn hình.

Bài toán đặt ra là cần một mô hình dữ liệu và mô hình giao diện đủ tổng quát để chứa nhiều loại tương tác, nhưng vẫn đủ chặt chẽ để máy chủ có thể kiểm tra, chấm điểm và lưu kết quả một cách an toàn. Ngoài ra, mỗi loại tương tác phải có đủ ba góc nhìn: giảng viên biên tập, học viên trả lời và giảng viên/học viên xem lại kết quả.

\subsection{Giải pháp}
Giải pháp được sử dụng là mô hình hóa hoạt động tương tác theo kiểu ``một lõi chung, nhiều cấu hình chuyên biệt''. Mỗi hoạt động đều có các thuộc tính chung như định danh bài giảng, định danh phân đoạn, loại tương tác, mốc thời gian xuất hiện, thứ tự, điểm tối đa, nội dung câu hỏi và giải thích. Phần khác nhau giữa các loại tương tác được đặt trong cấu hình riêng theo từng loại, ví dụ cấu hình trắc nghiệm, điền khuyết, nghe hiểu, đọc hiểu và viết.

Ở tầng giao tiếp, hợp đồng Protocol Buffers định nghĩa rõ các loại cấu hình và phản hồi tương ứng. Cách định nghĩa này giúp máy chủ và giao diện dùng chung một kiểu dữ liệu, giảm rủi ro sai lệch khi truyền nhận qua Connect RPC. Ở tầng máy chủ, mỗi loại phản hồi được chấm theo chiến lược phù hợp: trắc nghiệm và điền khuyết có thể chấm bằng luật xác định; đọc hiểu, phát âm hoặc viết có thể cần gọi thêm dịch vụ trí tuệ nhân tạo để tạo điểm và phản hồi bằng ngôn ngữ tự nhiên. Hệ thống cũng hỗ trợ chế độ phản hồi khác nhau, chẳng hạn phản hồi sau từng câu hoặc phản hồi sau khi nộp toàn bộ bài làm.

Ở tầng giao diện, hệ thống sử dụng cơ chế sổ đăng ký renderer cho các loại tương tác. Mỗi loại tương tác cung cấp bộ thành phần riêng cho ba màn hình: màn hình làm bài của học viên, màn hình biên tập của giảng viên và dòng xem lại kết quả. Khi cần thêm một loại tương tác mới, lập trình viên không phải sửa toàn bộ trình phát bài học, mà chỉ cần bổ sung cấu hình dữ liệu, logic chấm điểm và renderer tương ứng rồi đăng ký vào sổ đăng ký chung. Thiết kế này giúp phần trình phát video và luồng nộp bài không phụ thuộc trực tiếp vào từng loại bài tập cụ thể.

Một điểm đáng chú ý là hệ thống không chỉ hỗ trợ các câu hỏi sinh tự động. Giảng viên có thể tạo thủ công hoặc chỉnh sửa hoạt động đã sinh. Nhờ dùng cùng một mô hình dữ liệu cho hoạt động sinh bởi AI và hoạt động do giảng viên tạo, hệ thống tránh phải duy trì hai đường xử lý riêng biệt. Thuộc tính nguồn tạo giúp phân biệt hoạt động tự động và thủ công khi cần thống kê hoặc rà soát.

\subsection{Kết quả đạt được}
Hệ thống đã hỗ trợ nhiều dạng tương tác gồm trắc nghiệm một đáp án, trắc nghiệm nhiều đáp án, điền khuyết, nghe hiểu, đọc hiểu và viết. Các dạng này cùng xuất hiện trong một bài học, cùng được đặt trên dòng thời gian video và cùng được lưu kết quả trong một luồng nộp bài thống nhất. Giảng viên có thể biên tập từng dạng bài trong giao diện quản lý bài giảng; học viên có thể trả lời ngay khi video tạm dừng tại mốc tương tác; kết quả được chấm điểm, phản hồi và lưu lại để phục vụ phân tích.

Đóng góp của giải pháp là tạo ra một mô hình tương tác mở rộng được thay vì chỉ cài đặt một loại câu hỏi cố định. Điều này làm cho hệ thống phù hợp hơn với nhiều loại nội dung học tập, đồng thời giảm chi phí phát triển khi cần bổ sung dạng bài mới trong tương lai.

\section{Xử lý tác vụ dài bằng hàng đợi bền vững}
\subsection{Đặt vấn đề}
Quy trình phân tích video là một tác vụ dài, phụ thuộc vào nhiều dịch vụ bên ngoài và có thể kéo dài hơn nhiều so với một yêu cầu HTTP thông thường. Nếu xử lý toàn bộ trong lúc người dùng chờ phản hồi, giao diện dễ bị treo, kết nối có thể hết thời gian chờ và khi máy chủ khởi động lại thì kết quả đang xử lý bị mất. Ngoài ra, một video có thể thất bại ở một bước cụ thể như phiên âm hoặc sinh bài tập; nếu hệ thống không lưu trạng thái trung gian, người dùng chỉ biết là tác vụ thất bại mà không biết đang lỗi ở đâu và hệ thống cũng khó chạy lại đúng cách.

Do đó, hệ thống cần một cơ chế xử lý nền có khả năng lưu trạng thái bền vững, cho phép nhiều worker xử lý song song, chống tình trạng hai worker cùng ghi kết quả cho một tác vụ, đồng thời có khả năng phục hồi khi worker bị dừng giữa chừng. Đây là yêu cầu quan trọng vì pipeline AI không chỉ tốn tài nguyên mà còn tạo ra dữ liệu cần nhất quán với bài giảng và các hoạt động tương tác.

\subsection{Giải pháp}
Giải pháp được xây dựng là một hàng đợi tác vụ bền vững dựa trên PostgreSQL. Bảng tác vụ là nguồn dữ liệu chính về trạng thái xử lý, bao gồm loại tác vụ, dữ liệu đầu vào, trạng thái, tiến độ, thông điệp hiển thị, worker đang sở hữu và dữ liệu checkpoint. Khi giảng viên kích hoạt phân tích video, hệ thống không chạy pipeline ngay trong yêu cầu đó mà tạo một bản ghi tác vụ ở trạng thái chờ xử lý. Worker nền nhận tác vụ, chuyển trạng thái sang đang xử lý và cập nhật tiến độ sau mỗi giai đoạn.

Lớp hàng đợi được thiết kế độc lập với nghiệp vụ cụ thể. Nó không biết tác vụ ``phiên âm'' hay ``sinh bài tập'' có ý nghĩa gì; nó chỉ quản lý vòng đời tác vụ. Mỗi loại tác vụ cài đặt một giao diện executor riêng và được đăng ký vào hệ thống. Với pipeline phân tích video, executor điều phối ba bước: phiên âm, phân đoạn và sinh hoạt động tương tác. Sau khi hoàn thành một bước lớn, executor ghi checkpoint để biết bước nào đã xong. Nếu worker bị dừng và tác vụ được nhận lại, hệ thống có thể bỏ qua bước đã hoàn thành thay vì chạy lại toàn bộ từ đầu.

Để tránh lỗi khi nhiều worker cùng tham gia, mỗi worker có một định danh riêng. Khi worker nhận tác vụ, định danh này được ghi vào bản ghi tác vụ. Các thao tác cập nhật kết thúc tác vụ chỉ hợp lệ nếu bản ghi vẫn thuộc về worker đó và vẫn ở trạng thái đang xử lý. Nhờ vậy, nếu một worker cũ bị treo rồi quay lại sau khi tác vụ đã được worker khác nhận, worker cũ không thể ghi đè kết quả cuối. Cơ chế này phù hợp với mô hình phân phối ít nhất một lần: một tác vụ có thể được nhận lại, nhưng việc ghi trạng thái cuối được bảo vệ bằng điều kiện sở hữu.

Hàng đợi cũng hỗ trợ cập nhật tiến độ theo các nhãn cấp cao như phiên âm, phân đoạn và sinh bài tập. Giao diện có thể đọc các thông tin này để hiển thị tiến trình cho giảng viên thay vì chỉ hiển thị một trạng thái chung chung. Khi có lỗi, trạng thái lỗi và thông điệp lỗi được lưu lại để người dùng biết tác vụ dừng ở bước nào và có thể chạy lại sau khi điều chỉnh.

\subsection{Kết quả đạt được}
Nhờ cơ chế hàng đợi bền vững, yêu cầu tạo bài học từ video được phản hồi nhanh, còn phần xử lý nặng chạy nền một cách có kiểm soát. Hệ thống có thể theo dõi tiến độ từng tác vụ, phục hồi từ checkpoint khi tác vụ bị gián đoạn, và tránh ghi đè kết quả bởi worker không còn quyền sở hữu. Điều này giúp pipeline AI phù hợp hơn với môi trường triển khai thực tế, nơi container có thể khởi động lại, dịch vụ ngoài có thể chậm hoặc lỗi, và nhiều bài giảng có thể được phân tích trong cùng một thời điểm.

Đóng góp của phần này là thiết kế được một cơ chế xử lý nền tổng quát, không gắn cứng với riêng pipeline video. Cơ chế executor theo loại tác vụ cho phép hệ thống tiếp tục bổ sung các tác vụ dài khác trong tương lai, ví dụ sinh lại âm thanh, đồng bộ lại dữ liệu phân tích hoặc xử lý thống kê lớn.

\section{Phân tích kết quả học tập từ dữ liệu tương tác}
\subsection{Đặt vấn đề}
Một hạn chế phổ biến của các hệ thống phát video bài giảng là dữ liệu thu được thường chỉ dừng ở mức ``đã xem'' hoặc ``chưa xem''. Chỉ số này không đủ để giảng viên đánh giá học viên có hiểu bài hay không, cũng không chỉ ra phần kiến thức nào đang gây khó khăn cho cả lớp. Ngay cả khi có bài tập đi kèm, nếu kết quả bài làm không được gắn với mốc thời gian và phân đoạn nội dung, giảng viên vẫn khó biết câu hỏi sai thuộc phần nào của video và cần hỗ trợ ai.

Với Dyadia, bài toán không chỉ là chấm điểm từng câu trả lời, mà còn là biến dữ liệu tương tác thành thông tin có ích cho giảng viên. Hệ thống cần lưu được tiến độ xem video, kết quả từng hoạt động, điểm số, phản hồi và thời điểm làm bài; từ đó tổng hợp thành các chỉ báo như tỷ lệ hoàn thành, điểm trung bình, tỷ lệ trả lời đúng theo phân đoạn và danh sách học viên cần chú ý.

\subsection{Giải pháp}
Giải pháp của hệ thống là gắn dữ liệu học tập với hai trục chính: trục thời gian của video và trục phân đoạn nội dung. Mỗi hoạt động tương tác có mốc xuất hiện trong video và có thể liên kết với một phân đoạn bản chép lời. Khi học viên trả lời, hệ thống lưu câu trả lời, điểm số và phản hồi tương ứng. Nhờ vậy, mỗi kết quả không chỉ là một điểm số riêng lẻ mà còn có ngữ cảnh: học viên sai ở hoạt động nào, hoạt động đó thuộc đoạn nội dung nào và xuất hiện tại thời điểm nào trong video.

Đối với tiến độ xem video, hệ thống sử dụng FoundationDB để lưu bitmap theo từng cặp học viên--bài giảng. Mỗi bit biểu diễn một giây video đã được xem. Cách biểu diễn này tiết kiệm dung lượng hơn nhiều so với việc ghi từng khoảng xem thành nhiều bản ghi quan hệ, đồng thời cho phép tính nhanh tỷ lệ video đã xem bằng thao tác đếm bit. Hệ thống còn giới hạn tốc độ tăng của vùng đã xem dựa trên thời gian thực để hạn chế trường hợp trình duyệt báo giả rằng học viên đã xem toàn bộ video trong một thời gian quá ngắn.

Ở phía giảng viên, các dữ liệu này được tổng hợp thành giao diện theo dõi kết quả học tập. Bản đồ nhiệt hiển thị tỷ lệ trả lời đúng theo từng phân đoạn, giúp giảng viên nhận ra phần kiến thức cả lớp gặp khó khăn. Danh sách học viên cần chú ý dựa trên các chỉ báo như tiến độ thấp, điểm trung bình thấp hoặc mức tương tác hạn chế. Ngoài ra, hệ thống còn hiển thị các thống kê như số học viên hoạt động, tiến độ trung bình, điểm trung bình và biểu đồ phân bố kết quả, giúp giảng viên có cái nhìn tổng quan trước khi xem chi tiết từng học viên.

\subsection{Kết quả đạt được}
Kết quả đạt được là hệ thống không chỉ dừng ở việc phát video và chấm bài, mà đã hình thành được vòng phản hồi học tập. Học viên nhận được phản hồi sau khi trả lời hoạt động tương tác; giảng viên nhận được dữ liệu tổng hợp để điều chỉnh cách hỗ trợ lớp học. Các phân đoạn có tỷ lệ đúng thấp được làm nổi bật, giúp việc can thiệp có trọng tâm hơn so với việc chỉ nhìn vào điểm trung bình cuối cùng.

Đóng góp của giải pháp là kết hợp dữ liệu xem video với dữ liệu làm bài theo cùng một mô hình bài học tương tác. Nhờ đó, hệ thống tạo được cơ sở dữ liệu học tập có ý nghĩa sư phạm hơn: mỗi điểm số gắn với nội dung cụ thể, mỗi tiến độ xem có bằng chứng theo thời gian và mỗi thống kê đều có thể truy ngược về hoạt động hoặc phân đoạn tương ứng.


\end{document}
